\section{Methodological Findings}
\label{sec:method-findings}

\subsection{Serving-Path Sensitivity of Folded-Context Policies}
\label{sec:serving}

\begin{table}[t]
\centering
\small
\begin{tabular}{lcc}
\toprule
Checkpoint & Chat-templated serving & Training-style serving \\
\midrule
No-RL base & 0.147 & 0.160--0.167 \\
\foldgrpo{} step-50 & 0.027 & 0.247 \\
\foldgrpo{} step-100 & 0.093 & 0.267 \\
\bottomrule
\end{tabular}
\caption{The same checkpoints scored through a standard OpenAI-style chat endpoint versus the
training stack's token-level validation path. The RL'd policies collapse under chat
re-templating; the base model is unaffected.}
\label{tab:serving}
\end{table}

Our first evaluation attempt served checkpoints behind a standard \texttt{/chat/completions}
endpoint and produced catastrophically low scores for the RL'd\fredaauto{``RL'd'' is informal; prefer ``RL-trained'' (three occurrences in this section, incl.\ the Table 4 caption).} models while leaving the base
model unaffected (\cref{tab:serving}). The cause is that per-turn chat re-templating drops
prior-turn reasoning content that token-level continuation preserves; the RL'd policies come to
depend on it. Outputs remain locally coherent, but trajectories shorten from $\approx$17 to
$\approx$4 turns. The practical implication is that deployment serving must match training
serving for context-folding agents---a caveat absent from the paper. All scores in
\cref{sec:results} therefore use the training-style path for every arm.

\subsection{Per-Trajectory Contamination Screening}
\label{sec:contamination}

\begin{table}[t]
\centering
\small
\begin{tabular}{lcc}
\toprule
Cell & Contaminated & Clean re-measurement \\
\midrule
\foldgrpo{} step-100, $T{=}1.0$ & 0.147 & 0.282 \\
\foldgrpo{} step-50, $T{=}1.0$ & 0.090 & 0.243 \\
No-RL base, greedy & 0.000 & 0.167 \\
No-RL base, $T{=}1.0$ & 0.093 & 0.168 \\
\bottomrule
\end{tabular}
\caption{Evaluation cells that overlapped infrastructure outages of the shared search/judge
node, before and after re-measurement on isolated infrastructure.}
\label{tab:contamination}
\end{table}

Four evaluation cells overlapped outages of the shared search/judge node
(\cref{tab:contamination}). Two failed an aggregate screen (more than 50 tool-call errors per
run) and were discarded immediately; two passed it with \emph{sub-threshold} contamination and
silently lost 10--15 points. The latter pair briefly supported a plausible but false
conclusion---that \foldgrpo{}'s process rewards make the policy fragile to sampling
temperature---which scheduled cross-check runs on self-contained workers overturned: all arms
are temperature-robust. We conclude that contamination screening for tool-dependent agent
evaluations must operate per trajectory (e.g., discarding trajectories containing failed tool
calls) rather than per run, because aggregate error counts can pass while enough trajectories
are silently degraded to move headline numbers by several points.
